\originalchapter{8.333作业1：热力学}
\psource{Mehran Kardar，8.333 Fall 2013，PS1，截止2013年9月18日；与阶段复习1的同题另作对应。}
\begin{exercise}[原题1(a)--(d)]\label{p:333-1-1}
弹性丝位移$x$所需张力$J=ax-bT+cTx$，定$x$热容$C_x=A(x)T$。(a)由Maxwell关系求$S_x|_T$。(b)证明$A$不依$x$。(c)取$S(0,0)=S_0$求熵。(d)以$T,J$求定张力热容$C_J$，注明稳定区。
\end{exercise}
\begin{solution}
$\dd E=T\dd S+J\dd x$，$\dd F=-S\dd T+J\dd x$，所以$S_x=-J_T=b-cx$。而$S_T=C_x/T=A(x)$，混合偏导要求$A'(x)=\partial_T(b-cx)=0$。积分$S=S_0+AT+bx-cx^2/2$。定$J$时$x=(J+bT)/(a+cT)$，$\partial_Tx|_J=(ab-cJ)/(a+cT)^2$，故
\begin{equation}
C_J=AT+\frac{T(ab-cJ)^2}{(a+cT)^3}.
\end{equation}
局部机械稳定需$\partial_xJ=a+cT>0$，通常热稳定还需$A>0$。有限温经验方程不必适用于所有$T$，$T=0$积分常数是题给的数学延拓条件。
\end{solution}
\begin{exercise}[原题2(a)--(d)]\label{p:333-1-2}
硬芯气满足$p(V-Nb)=N\kb T$，固定$N$、$C_V$不随温度。(a)求涉及$S_V|_T$的Maxwell关系。(b)求$\dd E$并证明$E$不依$V$。(c)证明$\gamma=C_p/C_V=1+N\kb/C_V$。(d)证明绝热$p(V-Nb)^\gamma$恒定。
\end{exercise}
\begin{solution}
$S_V=p_T=N\kb/(V-Nb)$，$\dd E=C_V\dd T+[Tp_T-p]\dd V=C_V\dd T$，所以$E=C_VT+E_0(N)$。定压膨胀$V-Nb=N\kb T/p$，焓导数$C_p=C_V+N\kb$。熵$S=C_V\log T+N\kb\log(V-Nb)+S_0(N)$，绝热$T(V-Nb)^{N\kb/C_V}$恒定，乘状态方程即所求。要求$V>Nb,T>0$。
\end{solution}
\begin{exercise}[原题3(a)--(e)]\label{p:333-1-3}
零场超导与正常金属热容$C_s=V\alpha T^3,C_n=V(\beta T^3+\gamma T)$，体积变化忽略。(a)第三定律求两熵。(b)零场转变无潜热，求$T_c$。(c)$E_n(0)=E_0,E_s(0)=E_0-V\Delta$，求有限温能量。(d)相平衡求$\Delta$。(e)采用Gaussian磁功$\dd E=T\dd S+B\dd M+\mu\dd N$，$M_s=-VB/(4\pi),M_n=0$，求使超导失稳的$B_c(T)=B_0[1-T^2/T_c^2]$与$B_0$。
\end{exercise}
\begin{solution}
熵$S_s=V\alpha T^3/3,S_n=V(\beta T^3/3+\gamma T)$，无潜热要求$S_s(T_c)=S_n(T_c)$，得$T_c^2=3\gamma/(\alpha-\beta)$，需$\alpha>\beta,\gamma>0$。能量$E_s=E_0-V\Delta+V\alpha T^4/4$、$E_n=E_0+V\beta T^4/4+V\gamma T^2/2$。自由能差
\begin{equation}
F_n-F_s=V[\Delta+(\alpha-\beta)T^4/12-\gamma T^2/2].
\end{equation}
$T_c$相等给$\Delta=3\gamma^2/[4(\alpha-\beta)]=\gamma T_c^2/4$，所以差$V\Delta(1-T^2/T_c^2)^2$。磁场自然势$G=E-TS-BM$有$G_B=-M$，超导额外$VB^2/(8\pi)$、正常不变。平衡给$B_c^2=8\pi\Delta(1-T^2/T_c^2)^2$，故$B_0=\sqrt{8\pi\Delta}$，此形式只用于$T\le T_c$与题示热容模型。
\end{solution}
\begin{exercise}[原题4(a)--(e)]\label{p:333-1-4}
用光子气的无穷小Carnot循环，从相邻温度$T,T+\dd T$、压强$p,p+\dd p$与体积区间$V,V+\dd V$推黑体能量。(a)功$W$。(b)等温膨胀吸热$Q$。(c)用Carnot效率联系二者。(d)已知$p=AT^4,A=\pi^2\kb^4/[45(\hbar c)^3]$、$E(T,0)=0$求$E$。(e)求绝热线。
\end{exercise}
\begin{solution}
循环面积$W=\dd p\dd V$（取发动机方向），$Q=[(\partial_VE)_T+p]\dd V$，效率$\dd T/T$给$(\partial_VE)_T+p=T\dd p/\dd T$。代$p=AT^4$得$E=3AT^4V$。$S=(E+pV)/T=4AT^3V$，绝热$VT^3$恒定，亦$pV^{4/3}$恒定。原图的相邻曲线信息已由这些状态条件完整表述，不把示意图面积当有限循环精确公式。
\end{solution}
\begin{exercise}[原题5(a)--(b)]\label{p:333-1-5}
(a)两个初温不同的物体只通过热交换趋向共同温度，利用热流方向证明总熵增加。(b)可调体积气体置于恒$T,p$热库，热库熵变$-(\Delta E+p\Delta V)/T$，证明平衡极小$G=E+pV-TS$。
\end{exercise}
\begin{solution}
取$\delta Q>0$从较热物体1流到较冷2，均匀体的熵变化$\delta S=\delta Q(1/T_2-1/T_1)>0$，积分给正总增量；若物体内部有不可逆温差还可增加更多熵。(b)$\Delta S_{\rm tot}=\Delta S_{\rm gas}-(\Delta E+p\Delta V)/T=-\Delta G/T\ge0$，所以固定热库条件下自发过程$G$下降，稳定平衡取可达状态中最低$G$。这是组合体系熵原理，并非开放气体自身熵必增。
\end{solution}
\begin{exercise}[原题6(a)--(b)]\label{p:333-1-6}
两物体热容常数$C_1,C_2$、$T_1>T_2>0$。(a)直接接触且无外界交换，求最终温度与熵变。(b)用可逆Carnot发动机在二者间传热，求最终温度与输出功。
\end{exercise}
\begin{solution}
(a)能量守恒给$T_F=(C_1T_1+C_2T_2)/(C_1+C_2)$，$\Delta S=C_1\log(T_F/T_1)+C_2\log(T_F/T_2)\ge0$，正号由加权对数凹性。(b)总熵守恒给$T_R=T_1^{C_1/(C_1+C_2)}T_2^{C_2/(C_1+C_2)}$，输出$W=C_1T_1+C_2T_2-(C_1+C_2)T_R$。$T_R\le T_F$来自算术--几何均值；更多初始内能可作为功输出。
\end{solution}
\begin{exercise}[原题7]\label{p:333-1-7}
发动机与一组$T_{max}=T_1>\cdots>T_n=T_{min}$热源交换，吸入$Q_\alpha^+>0$、放出$Q_\beta^->0$，$W=Q^+-Q^-$。证明效率不超过仅用最高最低温的Carnot效率，说明严格小于的条件。
\end{exercise}
\begin{solution}
周期热源熵要求$\sum Q_\alpha^+/T_\alpha^+\le\sum Q_\beta^-/T_\beta^-$。左侧$\ge Q^+/T_{max}$，右侧$\le Q^-/T_{min}$，所以$\eta=1-Q^-/Q^+\le1-T_{min}/T_{max}$。非零热量在中间温度吸入或释放、或循环不可逆时不等号严格；若只在两端交换且可逆，可等号。原题“less”应按实际多温流量理解，不能排除经典Carnot等号。
\end{solution}
\begin{exercise}[原题8(a)--(c)，选做]\label{p:333-1-8}
近隔离稀薄气云经引力形成太阳和行星。(a)轨道和结构更有序为何不违反第二定律？(b)太阳核反应从质子形成较重元素是否因此降低总熵？(c)地球生命与智能的组织增长如何与第二定律相容？
\end{exercise}
\begin{solution}
(a)可见空间结构并非总微观熵。引力收缩把势能释放为运动、热与辐射，粒子与辐射的状态数可增长；自引力系统的长程非加性还使普通气体直觉的均匀平衡不直接适用。(b)核组成的某些自由度减少不决定总熵，必须把释能、光子及其他产物和传输一起计入，恒定粒子种类的简单混合熵不覆盖核反应。(c)地球接受高温太阳光并向较冷空间辐射，吸能与排能的熵不同，开放子系统可形成局部组织同时向环境排更多熵。第二律约束包括热源与辐射环境的总熵，不以“复杂程度”替代$S$。
\end{solution}
